\section{Introduction}\label{sec:introduction}

Let $k$ be a global function field: a finite extension of $\bbF_q(t)$.
Let $G$ be an absolutely almost simple simply connected algebraic group
over $k$. The group $G(k)$ is considered as an abstract group. Its normal
subgroups nevertheless reflect the topologies of its local groups.
For a place $v$ of $k$, write $k_v$ for the completion and set
\[
 A(G)=\{v:\rank_{k_v}G=0\},\qquad
 H_A=\prod_{v\in A(G)}G(k_v).
\]
The set $A(G)$ is finite. Each of its factors is compact, and $H_A$ carries
the product topology. Let $\delta_A:G(k)\longrightarrow H_A$ be the
diagonal homomorphism. The Margulis--Platonov conjecture asserts that
every noncentral normal subgroup of $G(k)$ is obtained by pulling back
an open normal subgroup of $H_A$.

\begin{theorem}\label{thm:main}
Let $k$ be a global function field and let $G$ be an absolutely almost
simple simply connected $k$-group. If
$N\triangleleft G(k)$ and $N\not\subseteq Z(G(k))$, then
\[
                  N=\delta_A^{-1}(W)
\]
for an open normal subgroup $W\triangleleft H_A$. This holds in every
positive characteristic, including when $G$ is anisotropic over $k$.
\end{theorem}

If $A(G)$ is empty, the product $H_A$ is trivial, and the theorem says that
$G(k)$ has no proper noncentral normal subgroup. In general it identifies
all noncentral normal subgroups using finitely many compact local groups.
In particular, every finite quotient of $G(k)$ is continuous for the
topology induced by these local factors: its finite-index kernel is
Zariski dense and therefore noncentral. The substantive issue is this particular
local description: finite index alone does not identify which local
topology detects a normal subgroup.

\subsection{Historical context}

The conjecture asks whether the compact groups at the anisotropic places
account for every noncentral normal subgroup of the rational-point group.
Its origins lie in Kneser's simplicity question for quaternion norm-one
groups and Platonov's extension of that question to simply connected simple
groups. Margulis formulated the description by anisotropic local factors
over global fields in \cite[Russian original, Section~2.4.8]{Margulis1979}.
The finite-index assertion in his Corollary~2.4.9 is an essential first
step. The further problem is to show that each resulting finite quotient
comes from the specified local groups. This also explains the conjecture's
role in the congruence subgroup problem, which studies finite-index
subgroups of arithmetic groups and needs additional information about
their congruence completions
\cite[pp.~73--76]{RapinchukCSP1992}.

For anisotropic inner forms of type~$A$, the problem concerns norm-one
groups of division algebras. The quaternion case was completed by
Margulis in 1980, as recalled in Tomanov's historical account
\cite[p.~895]{Tomanov1992}. Tomanov later proved the conjecture for
division algebras of power-of-two index and reduced arbitrary index to
odd index \cite[Theorem and Corollaries~1--2, p.~896]{Tomanov1992}.
Platonov--Rapinchuk and Raghunathan developed
the arithmetic control of commutator subgroups
\cite{PR1984,Raghunathan1988}; its global-field form is used explicitly in
\cite[proof of Theorem~6.4]{RapinchukSegevSeitz2002}.
Rapinchuk--Potapchik reduced the number-field problem to excluding
nonabelian finite simple quotients of the multiplicative group of the
division algebra \cite[Theorem~2.1 and Corollary~2.5]{RapinchukPotapchik1996}.
Segev's division-algebra argument and Segev--Seitz's analysis of commuting
graphs supplied this exclusion
\cite[Theorem~A]{Segev1999}\cite[Theorems~1--3]{SegevSeitz2002}.
Rapinchuk--Segev--Seitz proved the stronger result that every finite quotient
of the multiplicative group of a finite-dimensional division algebra is
solvable; their Theorem~6.4 proves the Margulis--Platonov conjecture for
inner type~$A$ over arbitrary global fields
\cite{RapinchukSegevSeitz2002}.
Rapinchuk subsequently gave a shorter proof using two-generation of finite
simple groups \cite{Rapinchuk2006}.

Over a global function field, Harder's results restrict the anisotropic
absolutely almost simple simply connected groups to type~$A$
\cite{Harder}\cite[p.~914]{PrasadTriality}.
The isotropic case follows from the global Kneser--Tits theorem and
projective simplicity \cite[Theorem~8.1]{Gille}.
Together with the established inner-type-$A$ theorem, these results reduce
the present problem to anisotropic outer forms of type~$A$, represented by
special unitary groups. The earlier state of the conjecture is discussed
in \cite[Section~3.5]{PRsurvey}; general outer-type-$A$ applications were
still explicitly conditional on it in
\cite[Corollary~1.6]{RapinchukTralle2026}.
Theorem~\ref{thm:main} provides the required normal-subgroup description
for these groups in every characteristic.

\subsection{Relation to the number-field proof}

The immediate predecessor is the number-field theorem of
\cite{MPNF}. We use its approach through finite quotients of abstract power
subgroups, approximation within a prescribed coset, and the distinction
between finite characters and nonabelian simple quotients. More
specifically, the power-factor construction comes from
\cite[Proposition~2.6]{MPNF}, and the finite lattice calculation for the
symmetric norm torus is the argument of
\cite[Lemmas~2.7--2.9]{MPNF}.
The difference-function and exact-transitivity argument for finite
characters adapts \cite[Section~4]{MPNF}.
The commuting approximation and circle-extension reduction follow
\cite[Section~3.1 and Sections~3.3--3.7]{MPNF}; the invariant probability
law follows \cite[Section~5.2]{MPNF}.
The finite-group obstruction is \cite[Lemma~5.2 and Section~6]{MPNF}, and
the unitary induction follows \cite[Section~7]{MPNF}.
We reproduce the lattice calculation and the finite-group argument. We
also verify an additional subgroup-complement property of the chosen
finite-group subsets, which is needed for the present recurrence proof.

Several independent inputs retain their own roles. We use Prasad's
strong approximation theorem over function fields \cite{PrasadSA}; the local and global
cohomological results of Bruhat--Tits and Harder provide the structural
reductions \cite{BruhatTitsIII,Harder}.
Demarche--Harari's duality and local-global theorems supply the
positive-characteristic arithmetic for tori and homogeneous spaces
\cite{DemarcheHarariDuality,DemarcheHarariHomogeneous}.
The passage from a simple quotient to an action of the full unitary group
uses Feit--Seitz's theorem that class-preserving automorphisms of a finite
simple group are inner \cite[Theorem~C]{FeitSeitz}.
The circle-extension argument uses Prasad--Rapinchuk's metaplectic-kernel
computation \cite{PRmetaplectic}, while the invariant probability law uses
Howe--Moore decay in its nonarchimedean form
\cite{HoweMoore}\cite[Theorem~4.19 and Definition~4.23]{Ciobotaru}.

The changes in positive characteristic occur at specific points of this
argument. When the exponent is divisible by the characteristic, the power
map need not be locally invertible. We adapt the inherited construction
by products of powers to this setting, separating the inseparable part
of the exponent and using the openness of the relevant local power
subgroups. The
additive group of the field has finite exponent, so the integer-dilation
recurrence used in the number-field proof does not give the required
invariance. That proof uses the density theorem of
Furstenberg--Katznelson \cite[Theorem~B]{FurstenbergKatznelson} in
\cite[Section~5.5]{MPNF}. Here we prove a replacement by finite additive
Fourier analysis. Odd characteristic uses orthogonality on finite additive
subgroups and a quadratic parametrization; characteristic two uses additive
polynomials and finite-dimensional coordinates over power subfields.
The resulting invariance is then combined with the finite-group
obstruction and the unitary induction.

\subsection{The argument}

The structural reductions leave an anisotropic special unitary group
$S=\SU(D,*)$, where $D$ is a central simple algebra of degree $n\geq3$
over a separable quadratic extension $L/k$, and $*$ is a unitary
involution. Write $U=\U(D,*)$ for the full unitary group. Every noncentral
normal subgroup of $S(k)$ has finite index. It therefore contains the
subgroup $R$ generated by all $e$th powers for some positive integer $e$.
For this $e$, define
\[
 P_0=\overline{\delta_A(R)},\qquad
 V_0=\delta_A^{-1}(P_0),\qquad V=V_0/R.
\]
The closure is taken in the product of the anisotropic local groups.
The group $V$ is finite; it measures the possible difference between an
abstract power subgroup and the subgroup prescribed by its local closure.
We prove $V=1$ for every $e$ by induction on $n$, simultaneously over all
global function fields.

Two arithmetic constructions permit us to work with this finite group.
First, products of conjugated torus powers let us choose a rational point
in an \emph{actual prescribed coset} of $R$, while imposing conditions at
finitely many places and controlling the remaining places by a
codimension-two exclusion. A line in the parameter space avoids that
exclusion on reduction. The anisotropy of the torus makes the parameter
map invariant under scalar multiplication; this recovers approximation
at the place omitted from strong approximation. Second, a particular norm
torus has both the Hasse principle and weak approximation. This turns
locally solvable simultaneous norm conditions into exact rational
solutions. Sections~\ref{sec:approximation} and~\ref{sec:norm-tori}
establish these constructions.

Suppose first that $D$ is a division algebra and $n$ is odd. Commuting
approximation makes the action of $U(k)$ on $V$ preserve conjugacy
classes. If $V\ne1$, a simple quotient and a circle-extension argument
then produce one of two obstructions: a finite abelian character of
$U(k)$ nontrivial on $S(k)$, or a homomorphism from $U(k)$ onto a finite
nonabelian simple group whose restriction to $V_0$ is surjective and
which kills $R$. Section~\ref{sec:quotients} proves this dichotomy.

The simple quotient gives a finite coloring of the Hermitian elements
of $D$ through a Cayley parametrization. Local mixing supplies an invariant
probability law for all translated colorings, with uniform color
marginals. A finite-group argument selects a nonempty proper conjugacy-invariant
subset of the target group and constrains returns between colors inside
and outside that subset under fractional transformations. Additive
recurrence in characteristic $p$ forces the indicator of membership in
this subset to be translation invariant. The fractional transformations
then generate left rotations through every target color, forcing the
chosen subset to be either empty or the whole group. Sections~
\ref{sec:palettes} and~\ref{sec:recurrence} develop this contradiction;
Appendix~\ref{sec:finite-groups} proves the finite-group input. Finite
abelian characters require a different argument. Exact transitivity on
nonzero pairs in $D^2$ and the additive span of the full unitary group reduce them to the known inner-type-$A$ theorem in
Section~\ref{sec:characters}.

It remains to pass from odd division degree to all degrees.
Section~\ref{sec:induction} does this using smaller unitary groups and
quaternion norm-one groups. In even degree, a fixed quadratic Hermitian
subfield supplies a centralizer of half the degree. A factorization
separates an element into a quaternion factor and a factor in that
centralizer. The approximation construction controls the anisotropic
places of both factors, and the norm-torus theorem makes the factorization
rational. Degree four requires a separate two-norm correction. Uniform
control of local power subgroups allows this induction even when
$p$ divides $e$. Finally, Section~\ref{sec:conclusion} converts $V=1$ into
Theorem~\ref{thm:main}.
